\subsection{Stratified Random Sampling}
% Source integration: supplied Chapter 7 screenshots, Sections 7.5--7.6 and
% selected stratified-sampling problems.

Stratified sampling partitions the population into subpopulations and samples
within each one.  It can guarantee representation of important groups and can
reduce variance when strata are internally homogeneous or have different means.

\begin{definition} \label{def:stratified-sampling-notation} \hlt{Strata and the Stratified Mean}\\
Let $U$ be a finite population, $L\geq1$, nonempty subsets $U_1,\ldots,U_L\subseteq U$, $N_l\defeq|U_l|$, $N\defeq\sum_{l=1}^{L}N_l$, and fixed values $x_{il}\in\R$ for $i=1,\ldots,N_l$, $l=1,\ldots,L$.\\
For stratum samples of sizes $n_l$, write sampled values as $X_{il}$, $i=1,\ldots,n_l$.
\begin{enumerate}[label=(\roman*)]
\setlength{\itemsep}{0pt}
\item \hlt{Strata}: $U=\bigsqcup_{l=1}^{L}U_l$.
\item \hlt{Stratum weight}: $W_l\defeq N_l/N$.
\item \hlt{Stratum mean}: $\mu_l\defeq N_l^{-1}\sum_{i=1}^{N_l}x_{il}$.
\item \hlt{Stratum variance}: $\sigma_l^2\defeq N_l^{-1}\sum_{i=1}^{N_l}(x_{il}-\mu_l)^2$.
\item \hlt{Stratum sample mean}: $\overline{X}_l\defeq n_l^{-1}\sum_{i=1}^{n_l}X_{il}$.
\item \hlt{Stratified mean}: $\overline{X}_s\defeq\sum_{l=1}^{L}W_l\overline{X}_l$.
\end{enumerate}
\hlt{Structural Commentary}: With $\mu\defeq N^{-1}\sum_{l=1}^{L}\sum_{i=1}^{N_l}x_{il}$, $\mu=\sum_{l=1}^{L}W_l\mu_l$.
\end{definition}

\begin{theorem} \label{thm:stratified-mean-unbiased-variance} \hlt{Mean and Variance of the Stratified Estimate}\\
Let $L\geq1$, $N_l>1$, $1\leq n_l\leq N_l$, $N\defeq\sum_{l=1}^{L}N_l$, and $W_l\defeq N_l/N$.  Fix $x_{il}\in\R$ for $i=1,\ldots,N_l$, $l=1,\ldots,L$, and set $\mu_l\defeq N_l^{-1}\sum_{i=1}^{N_l}x_{il}$, $\sigma_l^2\defeq N_l^{-1}\sum_{i=1}^{N_l}(x_{il}-\mu_l)^2$, and $\mu\defeq\sum_{l=1}^{L}W_l\mu_l$.  Assume independently $S_l\sim\operatorname{Unif}\{s\subseteq\{1,\ldots,N_l\}:|s|=n_l\}$ for $l=1,\ldots,L$.  For random orderings $(J_{il})_{i=1}^{n_l}$ of $S_l$, set $X_{il}\defeq x_{J_{il},l}$, $\overline{X}_l\defeq n_l^{-1}\sum_{i=1}^{n_l}X_{il}$, and $\overline{X}_s\defeq\sum_{l=1}^{L}W_l\overline{X}_l$.  Then:
\begin{enumerate}[label=(\roman*)]
\setlength{\itemsep}{0pt}
\item \hlt{Mean}: the stratified mean is design-unbiased for the population mean, $\mathbb{E}\overline{X}_s=\mu$.
\item \hlt{Variance}: the exact design variance under independent within-stratum sampling is
\[
\operatorname{Var}(\overline{X}_s)
=\sum_{l=1}^{L}
W_l^2\frac{\sigma_l^2}{n_l}
\left(1-\frac{n_l-1}{N_l-1}\right).
\]
\item \hlt{Small fractions}: when all stratum sampling fractions are small, the variance is approximated by
\[
\operatorname{Var}(\overline{X}_s)
\approx \sum_{l=1}^{L}W_l^2\frac{\sigma_l^2}{n_l}.
\]
\end{enumerate}
\begin{proof}
\begin{enumerate}[label=(\roman*)]
\setlength{\itemsep}{0pt}
\item \hlt{Mean}: Within stratum $l$, $\mathbb{P}(J_{il}=a)=1/N_l$ for every $a\in\{1,\ldots,N_l\}$, so $\mathbb{E}X_{il}=\mu_l$ and $\mathbb{E}\overline{X}_l=\mu_l$.  Therefore $\mathbb{E}\overline{X}_s=\sum_{l=1}^{L}W_l\mathbb{E}\overline{X}_l=\sum_{l=1}^{L}W_l\mu_l=\mu$.
\item \hlt{Variance}: Independence across strata gives
\[
\operatorname{Var}(\overline{X}_s)
=\sum_{l=1}^{L}W_l^2\operatorname{Var}(\overline{X}_l),
\]
and within stratum $l$ each sampled position has variance $\sigma_l^2$.  If $(Y_{1,l},\ldots,Y_{N_l,l})$ is a uniformly random ordering of $\{x_{il}:1\leq i\leq N_l\}$ and $\gamma_l$ is the common covariance of $Y_{a,l},Y_{b,l}$ for $a\neq b$, then $0=\operatorname{Var}(\sum_{i=1}^{N_l}Y_{i,l})=N_l\sigma_l^2+N_l(N_l-1)\gamma_l$, so $\gamma_l=-\sigma_l^2/(N_l-1)$ and $\operatorname{Var}(\overline{X}_l)=\sigma_l^2n_l^{-1}(1-(n_l-1)/(N_l-1))$.
\item \hlt{Small fractions}: If all stratum sampling fractions are small, each finite population correction is close to $1$.
\end{enumerate}
\hlt{Structural Commentary}: Stratified estimation is a weighted sum of independent within-stratum estimators.\\
\hlt{Key Insight}: Stratification controls variance by isolating within-stratum variation before averaging across the population weights.
\end{proof}
\end{theorem}

\begin{corollary} \label{cor:stratified-total-and-estimated-variance} \hlt{Totals and Estimated Variances}\\
Let $L\geq1$, $N_l>1$, $1<n_l\leq N_l$, $N\defeq\sum_{l=1}^{L}N_l$, and $W_l\defeq N_l/N$.  Fix $x_{il}\in\R$ for $i=1,\ldots,N_l$, $l=1,\ldots,L$, and set $\mu_l\defeq N_l^{-1}\sum_{i=1}^{N_l}x_{il}$, $\mu\defeq\sum_{l=1}^{L}W_l\mu_l$, and $\tau\defeq N\mu$.  Assume independently $S_l\sim\operatorname{Unif}\{s\subseteq\{1,\ldots,N_l\}:|s|=n_l\}$ for $l=1,\ldots,L$.  For random orderings $(J_{il})_{i=1}^{n_l}$ of $S_l$, set $X_{il}\defeq x_{J_{il},l}$, $\overline{X}_l\defeq n_l^{-1}\sum_{i=1}^{n_l}X_{il}$, $\overline{X}_s\defeq\sum_{l=1}^{L}W_l\overline{X}_l$, $T_s\defeq N\overline{X}_s$, and $s_l^2\defeq (n_l-1)^{-1}\sum_{i=1}^{n_l}(X_{il}-\overline{X}_l)^2$.  If $0<\alpha<1$ and $z(\alpha/2)\defeq\Phi^{-1}(1-\alpha/2)$ for the standard normal cdf $\Phi$, then:
\begin{enumerate}[label=(\roman*)]
\setlength{\itemsep}{0pt}
\item \hlt{Total moments}: the stratified total estimator is design-unbiased and scales the stratified mean variance, $\mathbb{E}T_s=\tau$ and $\operatorname{Var}(T_s)=N^2\operatorname{Var}(\overline{X}_s)$.
\item \hlt{Estimated mean variance}: the plug-in estimated variance of the stratified mean is
\[
s_{\overline{X}_s}^2
=\sum_{l=1}^{L}W_l^2\frac{s_l^2}{n_l}
\left(1-\frac{n_l}{N_l}\right).
\]
\item \hlt{Intervals}: with standard normal right-tail quantile $z(\alpha/2)$, approximate confidence intervals for $\mu$ and $\tau$ are $\overline{X}_s\pm z(\alpha/2)s_{\overline{X}_s}$ and $T_s\pm z(\alpha/2)Ns_{\overline{X}_s}$.
\end{enumerate}
\hlt{Structural Commentary}: Total estimation and interval construction are obtained by scaling the stratified mean estimate.\\
\hlt{Key Insight}: The variance estimate is additive across independent strata after each stratum is weighted by $W_l^2$.
\end{corollary}

\begin{theorem} \label{thm:neyman-allocation} \hlt{Neyman Allocation}\\
Let $L\geq1$, $W_l>0$ with $\sum_{l=1}^{L}W_l=1$, $\sigma_l>0$ for $l=1,\ldots,L$, and $n>0$.  For continuous allocations $n_l>0$ with $\sum_{l=1}^{L}n_l=n$, set
\[
V(n_1,\ldots,n_L)\defeq \sum_{l=1}^{L}W_l^2\frac{\sigma_l^2}{n_l}.
\]
The unique minimizer of $V$ is
\[
n_l
=n\frac{W_l\sigma_l}{\sum_{k=1}^{L}W_k\sigma_k},
\]
for $l=1,\ldots,L$.
\begin{proof}
The objective is convex in the variables $n_l>0$, and its Lagrangian is
\[
\mathcal{L}(n_1,\ldots,n_L,\lambda)
=\sum_{l=1}^{L}\frac{W_l^2\sigma_l^2}{n_l}
+\lambda\left(\sum_{l=1}^{L}n_l-n\right).
\]
The first-order equations give $-W_l^2\sigma_l^2/n_l^2+\lambda=0$, hence $n_l=W_l\sigma_l/\sqrt{\lambda}$.  Summing over $l$ and using $\sum_{l=1}^{L}n_l=n$ determines $1/\sqrt{\lambda}=n/\sum_{k=1}^{L}W_k\sigma_k$.
\hlt{Structural Commentary}: The constrained optimum balances marginal variance reduction across strata.\\
\hlt{Key Insight}: A stratum receives more sample not only when it is large, but also when it is internally variable.
\end{proof}
\end{theorem}

\begin{corollary} \label{cor:stratified-allocation-comparisons} \hlt{Allocation Comparisons}\\
Let $L\geq1$, $W_l>0$ with $\sum_{l=1}^{L}W_l=1$, $\mu_l\in\R$, $\sigma_l>0$, $\mu\defeq\sum_{l=1}^{L}W_l\mu_l$, $\overline{\sigma}\defeq\sum_{l=1}^{L}W_l\sigma_l$, and $n>0$.  Ignoring finite population corrections, let $\overline{X}_{so}$ denote the optimal-allocation stratified mean, let $\overline{X}_{sp}$ denote the proportional-allocation stratified mean with $n_l=nW_l$, and let $\overline{X}$ denote the ordinary simple-random-sampling mean with sample size $n$.  Then:
\begin{enumerate}[label=(\roman*)]
\setlength{\itemsep}{0pt}
\item \hlt{Allocation variances}: the approximate variances under optimal and proportional allocation are
\[
\operatorname{Var}(\overline{X}_{so})
=\frac{\left(\sum_{l=1}^{L}W_l\sigma_l\right)^2}{n},
\;
\operatorname{Var}(\overline{X}_{sp})
=\frac1n\sum_{l=1}^{L}W_l\sigma_l^2.
\]
\item \hlt{Optimal-allocation gain}: the approximate variance reduction from optimal allocation relative to proportional allocation is
\[
\operatorname{Var}(\overline{X}_{sp})
-\operatorname{Var}(\overline{X}_{so})
=\frac1n\sum_{l=1}^{L}W_l(\sigma_l-\overline{\sigma})^2.
\]
\item \hlt{Proportional-allocation gain}: the approximate variance reduction from proportional stratification relative to ordinary SRS is
\[
\operatorname{Var}(\overline{X})
-\operatorname{Var}(\overline{X}_{sp})
=\frac1n\sum_{l=1}^{L}W_l(\mu_l-\mu)^2.
\]
\end{enumerate}
\hlt{Structural Commentary}: Proportional allocation removes between-stratum mean variation, and optimal allocation additionally responds to unequal within-stratum standard deviations.\\
\hlt{Key Insight}: Stratification helps when means differ across strata; Neyman allocation helps further when precisions differ across strata.
\end{corollary}

\begin{example} \label{ex:hospital-stratification-summary} \hlt{Hospital Stratification by Beds}\\
In the hospital example, strata formed by number of beds have the following
population summaries:
\[
\begin{array}{c|c|c|c|c}
\text{Stratum} & N_l & W_l & \mu_l & \sigma_l\\ \hline
A & 98 & .249 & 182.9 & 103.4\\
B & 98 & .249 & 526.5 & 204.8\\
C & 98 & .249 & 956.3 & 243.5\\
D & 99 & .251 & 1591.2 & 419.2
\end{array}
\]
Equal stratum sample sizes give, ignoring finite population corrections,
\[
\operatorname{Var}(\overline{X}_s)=\frac{72042.6}{n},
\;
\sigma_{\overline{X}_s}=\frac{268.4}{\sqrt n}.
\]
The comparable simple-random-sampling standard deviation is about
$587.7/\sqrt n$, so stratifying by bed count substantially reduces variance.
For Neyman allocation, the allocation weights
$W_l\sigma_l/\sum_kW_k\sigma_k$ are approximately
\[
.106,\; .210,\; .250,\; .434,
\]
so the largest, most variable stratum receives the heaviest sampling effort.
\end{example}

\begin{remark} \label{rem:survey-sampling-model-limitations} \hlt{Model Limits and Other Designs}\\
Systematic sampling and cluster sampling are common designs beyond this first
development.  Systematic sampling can resemble simple random sampling when the
list order is effectively random, but periodic structure may create bias.
Cluster sampling is often cheaper, but it is less precise than simple random
sampling when units inside a cluster are similar.  More broadly, the standard
errors above measure random sampling variation, not nonresponse, selection
bias, interviewer effects, or question wording effects.
\end{remark}
